% !TeX root = ../../Book3.tex
% 纲要标题：### 7.3 纤维与基域扩张
% 纲要位置：工作记录/计划纲要.md，第 2020 行。
% 纲要细目（写作范围）：
% * 纤维的方程
% * 张量积和坐标环
% * 非代数闭域上的点与几何点

\section{纤维与基域扩张}
\label{b3:sec:0703}

把参数固定为一个值，就得到映射的一条纤维。这个操作在方程中是代入，在代数中是取商或张量积。取商可能产生幂零元，即使原来的源和目标都约化；因此纤维的点集与记录全部关系的代数必须分别保留。

\subsection{纤维的方程}
\label{b3:subsec:0703:01}

设 \(k\) 仍代数闭，\(\varphi:X\to Y\) 为态射，\(A=k[X],B=k[Y]\)。对 \(b\in Y\)，记其极大理想为 \(\mathfrak m_b\)，并令 \(J=\varphi^*(\mathfrak m_b)A\)。若坐标表达为 \(\varphi=(f_1,\ldots,f_m)\)、\(b=(b_1,\ldots,b_m)\)，则
\[
\varphi^{-1}(b)=Z_X(f_1-b_1,\ldots,f_m-b_m)=Z_X(J).
\]
这个闭集称为 \(b\) 上的\term[xianwei]{纤维}[fiber]。用强零点定理得到它作为代数集的坐标环
\begin{equation}\label{b3:eq:reduced-fiber-ring}
k\bigl[\varphi^{-1}(b)\bigr]\cong A/\sqrt J.
\end{equation}

\begin{proposition}\label{b3:prop:fiber-algebra}
设 \(k\) 为代数闭域，\(\varphi:X\to Y\) 为 \(k\) 上仿射代数集之间的态射，\(A=k[X]\)、\(B=k[Y]\)，\(b\in Y\)。令 \(\mathfrak m_b=\{g\in B:g(b)=0\}\)，\(J=\varphi^*(\mathfrak m_b)A\)。则有自然同构
\[
A\otimes_B\kappa(b)\cong A/J,
\qquad \kappa(b)=B/\mathfrak m_b\cong k.
\]
因此纤维的坐标环是该张量代数的约化商；只有当 \(J\) 已为根理想时，才能直接把坐标环写成 \(A/J\)。
\end{proposition}
\begin{proof}
将 \(a\otimes(c+\mathfrak m_b)\) 送到 \(a\varphi^*(c)+J\)。改变 \(c\) 的代表所产生的差在 \(J\) 中，且该公式对 \(B\) 平衡、保持代数乘法，故得到同态。反向由 \(a+J\mapsto a\otimes1\) 给出；\(J\) 的每个生成元都在这个映射下为零，所以良定义。两个复合在纯张量及 \(a+J\) 上分别为恒等，得到同构。再结合\eqref{b3:eq:reduced-fiber-ring}即得最后断言。
\end{proof}

\begin{table}[htbp]
\centering
\caption[纤维代数与集合纤维的坐标环]{纤维代数与集合纤维的坐标环。设 \(k\) 代数闭，\(\varphi:X\to Y\) 为仿射代数集的态射，\(A=k[X]\)、\(B=k[Y]\)、\(b\in Y\)，且 \(J=\varphi^*(\mathfrak m_b)A\)。}
\label{b3:tab:fiber-algebra-and-reduction}
\begin{tabularx}{\textwidth}{@{}l>{\raggedright\arraybackslash}X>{\raggedright\arraybackslash}p{0.3\textwidth}@{}}
\toprule
对象 & 代数表达 & 平方映射 \(t\mapsto t^2\) 的零纤维 \\
\midrule
纤维代数 & \(A\otimes_B\kappa(b)\cong A/J\) & \(k[t]/(t^2)\) \\
集合纤维的坐标环 & \(A/\sqrt J\) & \(k[t]/(t)\cong k\) \\
\bottomrule
\end{tabularx}
\end{table}

在\cref{b3:tab:fiber-algebra-and-reduction}的零纤维代数中，\(\bar t\) 非零却平方为零；普通点集只有零，因而其坐标环不保留这个类。对同一个平方映射，若 \(\operatorname{char}k\ne2\)，在一处的纤维代数为 \(k[t]/(t^2-1)\cong k\times k\)，有两个不同点；在特征二时它变成 \(k[t]/\bigl((t-1)^2\bigr)\)，仍只有一个点，却保留非零平方零元。相同的代入过程可以产生不同的根数与幂零结构。

\subsection{张量积与坐标环}
\label{b3:subsec:0703:02}

经典仿射代数集的坐标环把在整个点集上为零的多项式看作零，因此总是约化的。纤维代数则要保留代入后的全部关系，已经可以出现 \(k[t]/(t^2)\) 这样的非约化环。为了保留这些幂零信息及其扩域后的变化，研究对象便从普通点集转为有限型代数，相应的点空间是代数的素谱。设 \(k\) 为任意域，\(A\) 为交换含幺有限型 \(k\) 代数，不要求约化。对域扩张 \(K/k\)，记 \(A_K=K\otimes_kA\)，称为将标量扩张到 \(K\) 的代数。

\begin{proposition}\label{b3:prop:base-change-presentation}
设 \(k\) 为域，\(K/k\) 为域扩张，\(I\) 为 \(k[x_1,\ldots,x_n]\) 的理想，\(A=k[x_1,\ldots,x_n]/I\)，其中 \(n\ge0\)。则
\[
K\otimes_kA\cong K[x_1,\ldots,x_n]/I K[x_1,\ldots,x_n].
\]
进一步，设 \(m\ge0\)，\(J\) 为 \(k[y_1,\ldots,y_m]\) 的理想，\(B=k[y_1,\ldots,y_m]/J\)。令
\[
R=k[x_1,\ldots,x_n,y_1,\ldots,y_m],
\]
并记 \(a_i=x_i+I\in A\)、\(b_j=y_j+J\in B\)。
则两组关系在 \(R\) 中的扩张理想 \(IR,JR\) 给出自然同构
\[
\begin{aligned}
A\otimes_kB&\cong R/(IR+JR),\\
a_i\otimes1&\longleftrightarrow x_i+(IR+JR),\\
1\otimes b_j&\longleftrightarrow y_j+(IR+JR).
\end{aligned}
\]
\end{proposition}
\begin{proof}
第一式中，将 \(c\otimes\bar f\) 送到 \(cf\) 的商类，得到良定义的代数同态。反向将 \(K\) 中的标量送到 \(c\otimes1\)，变量送到 \(1\otimes\bar x_i\)；原关系都为零，所以该赋值可下降到商。两复合固定所有标量和代数生成元，因此互逆。

第二项把 \(x_i\) 送到 \(a_i\otimes1\)，把 \(y_j\) 送到 \(1\otimes b_j\)。两组像元彼此交换，各自满足 \(I\) 与 \(J\) 中的关系，故得到 \(R/(IR+JR)\to A\otimes_kB\)。反向将 \((f+I)\otimes(g+J)\) 送到 \(fg+(IR+JR)\)，其中 \(f\in k[x_1,\ldots,x_n]\)、\(g\in k[y_1,\ldots,y_m]\)。这个公式对 \(k\) 平衡，改变两个代表元产生的差分别属于 \(IR\) 与 \(JR\)，因而良定义。两个复合固定两组生成元及 \(k\) 中的标量，得到所述同构。
\end{proof}

这个结论说明张量积怎样保存方程，却不自动保证所得环约化。若把一组方程解释为代数闭域中的普通零点集，强零点定理仍要求取其关系理想的根。在不代数闭的域上，仅看本域中的点还可能损失更多信息。

\ProseParagraphBreak
例如 \(A=\mathbb R[u]/(u^2+1)\cong\mathbb C\) 没有到 \(\mathbb R\) 的实代数同态，但
\[
\mathbb C\otimes_{\mathbb R}A
\cong\mathbb C[u]/(u^2+1)\cong\mathbb C\times\mathbb C,
\]
两个因子对应 \(u=i,-i\)。从环的角度，原对象并不为空；只是它没有实坐标点。不能把实零点集为空与原代数为零环混为一谈。

\subsection{非代数闭域上的点与几何点}
\label{b3:subsec:0703:03}

\begin{definition}\label{b3:def:field-valued-point}
设 \(k\) 为域，\(A\) 为交换含幺有限型 \(k\) 代数，\(K/k\) 为域扩张。若 \(\varphi:A\to K\) 为 \(k\) 代数同态，则称 \(\varphi\) 为一个 \(K\) 值点。若 \(K=k\)，则称它为 \(k\) 有理点；若 \(K\) 为代数闭域，则称它为以 \(K\) 为值的\term[jihe-dian]{几何点}[geometric point]。
\end{definition}

若 \(A=k[x_1,\ldots,x_n]/I\)，一个这样的同态恰由使全部 \(I\) 关系为零的一组 \(K\) 坐标确定。它的核在 \(\operatorname{Spec}A\) 中确定一个点，但未必是极大理想：例如自然包含 \(k[t]\to\overline{k(t)}\)、\(t\mapsto t\) 给出一个几何点，其在 \(\operatorname{Spec}k[t]\) 中的像为 \((0)\)，是全谱的一般点而非闭点。这里 \(t\) 的像是超越元；经典仿射直线的 \(k\) 坐标点 \(a\) 则对应极大理想 \((t-a)\)。如果只取固定代数闭包 \(\bar k\) 中的值，则由有限型性，坐标的像生成有限代数扩域，核为极大理想。

\ProseParagraphBreak
对任意 \(\mathfrak p\in\operatorname{Spec}A\)，点处的剩余域 \(\kappa(\mathfrak p)\) 都可嵌入某个代数闭域，从而给出位于它上面的几何点。这与只枚举 \(k\) 有理点不同。更一般地，给定交换含幺有限型 \(k\) 代数之间的同态 \(B\to A\) 及底点 \(\mathfrak q\in\operatorname{Spec}B\)，纤维代数 \(A\otimes_B\kappa(\mathfrak q)\) 就是在固定这个底点后，继续研究剩余域扩张中坐标的对象。

\begin{example}[变基后的非约化性]\label{b3:exm:inseparable-geometric-thickening}
设 \(p\) 为素数，\(s\) 为 \(\mathbb F_p\) 上的超越元，\(k=\mathbb F_p(s)\)。在 \(k\) 的一个代数闭包中取 \(\alpha\) 满足 \(\alpha^p=s\)，并令 \(L=k(\alpha)\)。则 \(L\) 为域，但
\[
L\otimes_kL\cong L[\varepsilon]/(\varepsilon^p).
\]
\end{example}
\begin{proof}
\(s\) 不是 \(k\) 中的 \(p\) 次幂：若 \((f/g)^p=s\)，把 \(f,g\in\mathbb F_p[s]\) 取互素，比较不可约因子 \(s\) 的指数便矛盾。因为 \(\alpha^p=s\)，其首一最小多项式 \(m_\alpha(T)\) 在 \(k[T]\) 中整除 \(T^p-s\)。在代数闭包中，\(T^p-s=(T-\alpha)^p\)，所以 \(m_\alpha\) 的全部根都为 \(\alpha\)，即 \(m_\alpha(T)=(T-\alpha)^r\)，其中 \(1\le r\le p\)。若 \(r<p\)，其 \(T^{r-1}\) 系数 \(-r\alpha\in k\) 会使 \(\alpha\in k\)，与上述结论矛盾。故 \(L\cong k[T]/(T^p-s)\)。应用\cref{b3:prop:base-change-presentation}后，在 \(L\) 中写 \(T^p-s=(T-\alpha)^p\)，令 \(\varepsilon=T-\alpha\) 即得所述同构。
\end{proof}

因此“约化”与“经过任意扩域仍约化”是不同要求。上例变基后仍只出现一个点，却有 \(p\) 次幂零结构。本章保留这个具体区别；一般几何约化性、可分性和平展性的关系将在后面的微分与平展章节中继续讨论。
